\documentclass[10pt]{article}
\usepackage[margin=0.84in]{geometry}
\usepackage[T1]{fontenc}
\usepackage{lmodern,microtype}
\usepackage{amsmath,amssymb,amsthm,mathtools}
\usepackage{graphicx}
\usepackage{xcolor}
\usepackage[colorlinks=true,allcolors=blue!55!black]{hyperref}
\setlength{\parindent}{0pt}
\setlength{\parskip}{0.55em}
\newtheorem{theorem}{Theorem}
\newcommand{\M}{M_2(\mathbb C)}
\newcommand{\Tr}{\operatorname{Tr}}
\newcommand{\cL}{\mathcal L}
\newcommand{\W}{\mathcal W}

\begin{document}

\begin{center}
{\Large\bfseries SRD detailed balance does not imply GNS detailed balance}\par
\medskip
{\large A primitive qubit Lindblad counterexample}\par
\medskip
{\small Source target: arXiv:1810.00906, Remark IV.5(2)}
\end{center}

\textbf{Status: candidate full counterexample; likely valid pending human review.}

\section*{Source question}

Yu Cao, Jianfeng Lu, and Yulong Lu, \emph{Gradient flow structure and
exponential decay of the sandwiched R\'enyi divergence for primitive Lindblad
equations with GNS-detailed balance}, arXiv:1810.00906; J. Math. Phys.
\textbf{60} (2019), 052202, prove
\[
 \text{GNS detailed balance}\Longrightarrow
 \text{SRD detailed balance}\Longrightarrow
 \text{KMS detailed balance}.
\]
Remark IV.5(2), on PDF page 23, asks whether the first implication reverses:
\begin{quote}
It is still unknown whether the SRD detailed balance condition is equivalent
to the GNS detailed balance condition or not.
\end{quote}

\begin{center}
\includegraphics[width=0.91\textwidth]{figures/open_problem_crop.png}
\end{center}

For a faithful invariant state $\sigma$, the source defines SRD detailed
balance by requiring the Heisenberg generator $\cL$ to be self-adjoint for
every inner product
\[
 \langle A,B\rangle_{\W_{\sigma,\alpha}}
   =\Tr\!\left(A^*\W_{\sigma,\alpha}(B)\right),
 \qquad \alpha>0.
\]
Equivalently,
$\W_{\sigma,\alpha}\cL\W_{\sigma,\alpha}^{-1}=\cL^\dagger$.

\section*{Proof intuition}

In a basis diagonalizing a qubit state, every SRD weight gives the \emph{same}
weight to the two oppositely oriented coherences $E_{01}$ and $E_{10}$.  GNS
does not: it weights them by the two different eigenvalues of the invariant
state.  We therefore add a Pauli-$X$ noise term that mixes $E_{01}$ and
$E_{10}$.  This mixing is symmetric for every SRD geometry, but it is not GNS
symmetric when the stationary state is nontracial.  Asymmetric up/down jumps
make that stationary state nontracial while preserving a primitive GKSL
generator.

\section*{Counterexample theorem}

Let $E_{ij}=|i\rangle\langle j|$ in the standard basis of $\mathbb C^2$ and
$X=E_{01}+E_{10}$.  For $a,b,c>0$, define
\[
 V_+=\sqrt a\,E_{10},\qquad
 V_-=\sqrt b\,E_{01},\qquad
 V_x=\sqrt c\,X
\]
and the Heisenberg generator
\begin{equation}\label{eq:gksl}
 \cL(A)=\sum_{V\in\{V_+,V_-,V_x\}}
 \left(V^*AV-\frac12\{V^*V,A\}\right).
\end{equation}

\begin{theorem}
If $a,b,c>0$ and $a\ne b$, then \eqref{eq:gksl} generates a primitive
quantum Markov semigroup which satisfies SRD detailed balance for every
$\alpha>0$ but does not satisfy GNS detailed balance.  In particular, SRD
detailed balance is not equivalent to GNS detailed balance, already on a
qubit.
\end{theorem}

\begin{proof}
Put
\[
 A_0=a+c,\qquad B_0=b+c,\qquad
 p_0=\frac{B_0}{A_0+B_0},\qquad
 p_1=\frac{A_0}{A_0+B_0},\qquad
 \sigma=p_0E_{00}+p_1E_{11}.
\]
The expression \eqref{eq:gksl} is in GKSL form and each summand annihilates
the identity, so it generates a completely positive unital semigroup.  Direct
multiplication gives
\begin{align}
 \cL(E_{00})&=-A_0E_{00}+B_0E_{11},
 &\cL(E_{11})&=A_0E_{00}-B_0E_{11},\label{eq:diag}\\
 \cL(E_{01})&=-\gamma E_{01}+cE_{10},
 &\cL(E_{10})&=cE_{01}-\gamma E_{10}.\label{eq:coh}
\end{align}
where $\gamma=(a+b)/2+c$.  Thus, in the ordered Hilbert--Schmidt orthonormal
basis $(E_{00},E_{11},E_{01},E_{10})$,
\begin{equation}\label{eq:matrix}
 [\cL]=
 \begin{pmatrix}
 -A_0&A_0&0&0\\
 B_0&-B_0&0&0\\
 0&0&-\gamma&c\\
 0&0&c&-\gamma
 \end{pmatrix}.
\end{equation}
Since $p_0A_0=p_1B_0$, equation \eqref{eq:matrix} shows that
$\cL^\dagger(\sigma)=0$.

We now verify SRD symmetry.  The source paper diagonalizes its weight operator
in the matrix-unit basis.  For every $\alpha>0$ there is a positive number
$q_\alpha$ such that
\begin{equation}\label{eq:srdweight}
 [\W_{\sigma,\alpha}]
   =\operatorname{diag}(p_0,p_1,q_\alpha,q_\alpha).
\end{equation}
Indeed, for $\alpha\ne1$ its two off-diagonal coefficients are both
\[
 q_\alpha=(\alpha-1)
 \frac{p_0^{1/\alpha}-p_1^{1/\alpha}}
 {p_1^{(1-\alpha)/\alpha}-p_0^{(1-\alpha)/\alpha}},
\]
which is symmetric under $p_0\leftrightarrow p_1$; at $\alpha=1$ their common
value is the logarithmic mean $(p_0-p_1)/(\log p_0-\log p_1)$.

Multiplying \eqref{eq:matrix} by \eqref{eq:srdweight}, one obtains
\[
 \W_{\sigma,\alpha}\cL=\cL^\dagger\W_{\sigma,\alpha}.
\]
On the diagonal block this is exactly the identity $p_0A_0=p_1B_0$;
on the coherence block it follows because the block in \eqref{eq:coh} is real
symmetric and both coordinates have weight $q_\alpha$.  Since
$\W_{\sigma,\alpha}$ is invertible, the defining SRD identity holds for every
$\alpha>0$.

For GNS detailed balance, the source uses
$\langle A,B\rangle_1=\Tr(\sigma A^*B)$.  Relative to our ordered basis its
weight operator is right multiplication by $\sigma$, hence
\[
 [R_\sigma]=\operatorname{diag}(p_0,p_1,p_1,p_0).
\]
The coherence-flip entries in \eqref{eq:matrix} would be GNS symmetric only if
$cp_1=cp_0$.  More explicitly,
\[
 \langle\cL(E_{01}),E_{10}\rangle_1=cp_0,
 \qquad
 \langle E_{01},\cL(E_{10})\rangle_1=cp_1.
\]
Because $c>0$ and $a\ne b$ imply $p_0\ne p_1$, these quantities differ.
Thus GNS detailed balance fails.

Finally, the four eigenvalues of \eqref{eq:matrix} are
\[
 0,\quad -(A_0+B_0),\quad -\frac{a+b}{2},\quad
 -\frac{a+b}{2}-2c.
\]
The zero eigenspace is precisely $\mathbb C I$, all other eigenvalues are
strictly negative, and $\sigma$ is faithful.  Hence the predual semigroup
converges to $\sigma\Tr(\,cdot\,)$ and is primitive.
\end{proof}

\section*{Smallest explicit instance}

Take $(a,b,c)=(1,3,1)$.  Then
\[
 \sigma=\begin{pmatrix}2/3&0\\0&1/3\end{pmatrix},\qquad
 [\cL]=
 \begin{pmatrix}
 -2&2&0&0\\4&-4&0&0\\0&0&-3&1\\0&0&1&-3
 \end{pmatrix},
\]
with spectrum $\{0,-2,-4,-6\}$.  This single concrete generator already
answers the source question negatively; the three-parameter family shows that
the failure is robust rather than isolated.

\section*{Verification}

The accompanying script constructs the three Lindblad operators, derives the
superoperator from \eqref{eq:gksl}, checks the invariant state and unitality,
tests the SRD intertwining identity at seven R\'enyi orders, confirms a nonzero
GNS residual, and checks the exact predicted spectrum.  For the explicit
instance it reports maximum sampled SRD residual
$3.14\times10^{-16}$ and GNS residual $0.4714045$ in Frobenius norm.  These
computations are supplementary; the proof above is exact.

\section*{Novelty check, limitations, and review request}

A bounded search on 17 August 2026 used the exact arXiv id, title, DOI, authors,
the exact phrase ``SRD detailed balance'', ``sandwiched R\'enyi detailed
balance'', and combinations with GNS/KMS/equivalence/Lindblad.  Local run
indexes and web/arXiv-oriented searches found the source paper, the first
author's thesis restating the same open remark, and papers citing the gradient
flow result, but no later answer and no occurrence of this qubit family.  The
novelty assessment is therefore positive but bounded, not exhaustive.

This packet disproves only the proposed equivalence.  It does not characterize
all SRD-detailed-balanced generators and does not close the broader gap between
necessary and sufficient conditions for arbitrary R\'enyi gradient-flow
structures.  Human review should focus on the convention relating weighted
self-adjointness to $\W\cL=\cL^\dagger\W$, the orientation of the two GNS
weights, and the standard equivalence between the displayed spectral
calculation and primitivity.

\section*{References}

\begin{enumerate}
\item Y. Cao, J. Lu, and Y. Lu, \emph{Gradient flow structure and exponential
decay of the sandwiched R\'enyi divergence for primitive Lindblad equations
with GNS-detailed balance}, J. Math. Phys. \textbf{60} (2019), 052202;
arXiv:1810.00906.
\item V. Gorini, A. Kossakowski, and E. C. G. Sudarshan,
\emph{Completely positive dynamical semigroups of $N$-level systems},
J. Math. Phys. \textbf{17} (1976), 821--825.
\item G. Lindblad, \emph{On the generators of quantum dynamical semigroups},
Commun. Math. Phys. \textbf{48} (1976), 119--130.
\end{enumerate}

\end{document}
